\section{Related Work}
\label{sec:related}

% \subsection{Visuomotor Policies and Predictive Control}

% Behavior cloning provides the standard supervised formulation for visuomotor policy learning, mapping observations and task context directly to demonstrated actions. ACT and Diffusion Policy
% \cite{zhao2023learning, chi2025diffusion} established action chunking and generative action prediction as effective control primitives. RT-1/RT-2, OpenVLA, Octo, and the
% $\pi$ family \cite{brohan2022rt, brohan2023rt, kim2024openvla, black2024pi_0, intelligence2025pi_} extended this
% paradigm through language conditioning, large-scale pretraining, and cross-task or cross-embodiment learning. Open X-Embodiment \cite{o2024open} further demonstrated the value of heterogeneous robot
% data.

% Predictive policies expose an additional interface by modeling future
% observations or latent states. Dreamer and DayDreamer
% \cite{hafner2019dream, hafner2023mastering, wu2023daydreamer} learn latent dynamics for control, while GR-1/GR-2 and RoboDreamer \cite{wu2024unleashing, cheang2024gr, zhou2024robodreamer} use future prediction for policy learning or planning. WorldEval and Ctrl-World \cite{li2025worldeval, guo2026ctrl} use rollouts to assess policy behavior.

% Action-conditioned world models provide a complementary perspective by
% predicting how an environment evolves under a supplied action sequence.
% A2World and $\tau_0$-WM \cite{huang2026learning, zhou2026tau_0} develop such models
% for simulation, candidate evaluation, or policy learning, while
% World-Coherent Decoding \cite{zhang2026worldcoherent} treats future--action pairs
% as hypotheses that can be selected and checked after execution. Together, these methods show that predictive models can support action generation, simulation, and evaluation. However, action-conditioned prediction and video--action policy generation are often developed as separate components, leaving open how explicit prediction–action interfaces affect component reuse across tasks.

% \subsection{Visuomotor Policies and Predictive Control}

% Behavior cloning provides the standard supervised formulation for visuomotor
% policy learning, mapping observations and task context directly to demonstrated
% actions. ACT and Diffusion Policy
% \cite{zhao2023learning, chi2025diffusion} established action chunking and
% generative action prediction as effective control primitives. RT-1/RT-2,
% OpenVLA, Octo, and the $\pi$ family
% \cite{brohan2022rt, brohan2023rt, kim2024openvla, black2024pi_0,
% intelligence2025pi_} extended this paradigm through language conditioning,
% large-scale pretraining, and cross-task or cross-embodiment learning.
% Open X-Embodiment~\cite{o2024open} further demonstrated the value of
% heterogeneous robot data.

% Predictive policies expose an additional interface by modeling future
% observations or latent states. Early visual-foresight methods combined
% action-conditioned video prediction with model-predictive control for robotic
% manipulation~\cite{finn2017deep,ebert2018visual}. Dreamer and DayDreamer
% \cite{hafner2019dream, hafner2023mastering, wu2023daydreamer} learn latent
% dynamics for control, while GR-1/GR-2 and RoboDreamer
% \cite{wu2024unleashing, cheang2024gr, zhou2024robodreamer} use future
% prediction for policy learning or planning. WorldEval and Ctrl-World
% \cite{li2025worldeval, guo2026ctrl} use predictive rollouts to assess policy
% behavior.

% Action-conditioned world models provide a complementary perspective by
% predicting how an environment evolves under a supplied action sequence.
% A2World and $\tau_0$-WM
% \cite{huang2026learning, zhou2026tau_0} develop such models for simulation,
% candidate evaluation, or policy learning. GPC~\cite{qi2026inference} trains an action-conditioned world model from expert
% demonstrations and exploratory rollouts, then uses predictive look-ahead to
% rank and refine policy proposals at inference time.
% World-Coherent Decoding~\cite{zhang2026worldcoherent} similarly treats
% future--action pairs as hypotheses that can be selected and checked after
% execution. Together, these methods show that predictive models can support
% action generation, simulation, planning, and evaluation. However,
% action-conditioned prediction and video--action policy generation are often
% developed as separate components, leaving open how explicit
% prediction--action interfaces affect component reuse across tasks.

\subsection{Visuomotor Policies and Predictive Control}

Behavior cloning provides the standard supervised formulation for visuomotor
policy learning, mapping observations and task context directly to demonstrated
actions. Action chunking and generative action prediction have become effective
control primitives~\cite{zhao2023learning, chi2025diffusion}, while large-scale
vision--language--action models extend this paradigm through language
conditioning, pretraining, and cross-task or cross-embodiment learning
\cite{brohan2022rt, brohan2023rt, kim2024openvla, black2024pi_0,
intelligence2025pi_, o2024open, team2024octo}.

Predictive control introduces an additional interface by explicitly modeling
future observations or latent states. Early visual-foresight methods combined
action-conditioned video prediction with model-predictive control for robotic
manipulation~\cite{finn2017deep,ebert2018visual}. Subsequent work learns latent
or visual dynamics to support policy learning, planning, and imagined rollouts
\cite{hafner2019dream, hafner2023mastering, wu2023daydreamer,
wu2024unleashing, cheang2024gr, zhou2024robodreamer}. Predictive rollouts have also been used to evaluate policy behavior~\cite{li2025worldeval, guo2026ctrl}.

Action-conditioned world models provide a complementary perspective by
predicting how an environment evolves under supplied or candidate actions.
Recent methods use such predictions for simulation, candidate evaluation,
proposal refinement, and self-verifying planning
\cite{huang2026learning, zhou2026tau_0, qi2026inference,
zhang2026worldcoherent}. Together, these lines show that predictive models can
support action generation, simulation, planning, and evaluation. However,
action-conditioned prediction and video--action policy generation are often
developed as separate components, leaving open how explicit
prediction--action interfaces affect component reuse across tasks.

% \subsection{Video Foundations and World--Action Architectures}

% Recent world-action models use pretrained video representations to connect
% future-observation prediction with robot control. DreamZero and VideoVLA
% \cite{ye2026world, shen2026videovla} jointly generate future video and action
% chunks, while LingBot-VA \cite{li2026causal} uses a causal video-first
% formulation. EgoWAM and Zero-WAM
% \cite{li2026egowam, zhou2026zero} use egocentric video as a source of predictive
% priors or task specifications, and ZimaBlue
% \cite{wu2026zimablue} studies multi-stage video pretraining followed by robot
% grounding.

% A separate line studies how predictive computation should be connected to
% action generation. Fast-WAM, World Tokens, and GlanceWAM
% \cite{yuan2026fast, tang2026world, wang2026glancewam} reduce,
% sparsify, or reschedule future prediction during inference.
% SelfWAM \cite{pan2026selfwam} conditions future prediction on clean actions during
% training while retaining an action-only deployment path.
% OA-WAM and RepWAM \cite{liu2026oa, wang2026repwam} change the predictive representation through object-addressable or action-grounded latents, while GIFT, SA-WAM, and TacPAC
% \cite{zheng2026gift, gonzalez2026spatially, ma2026tacpac} add structural, geometric, or tactile supervision.
% % OA-WAM~\cite{liu2026oa} explores object-addressable predictive representations, while TacPAC~\cite{ma2026tacpac} incorporates tactile supervision into world-action modeling.
% UVA~\cite{li2503unified} jointly trains
% policy, inverse-dynamics, and forward-dynamics objectives within
% a shared video--action model. 
% Riemann-1.0 and Motus2 \cite{sun2026riemann, bi2026motus2} expose policy and action-conditioned simulation within broader shared-model systems; Motus2 additionally includes
% outcome evaluation and policy-improvement loops.
% % Motus2~\cite{bi2026motus2} exposes policy and action-conditioned simulation within a shared-model system, while additionally supporting outcome evaluation and policy-improvement loops.
% AcrossWAM1.0 \cite{zhang2026acrosswam1} makes the boundaries between latent prediction, policy adaptation, and deployment explicit.

% These studies reveal several coupled design variables, including the
% prediction target, cross-modal information flow, generation order, and
% inference schedule. \ours fixes the causal robot-video initialization and
% training interfaces while comparing joint, sequential, and restricted
% video--action interaction under a shared evaluation protocol.

% \subsection{Video Foundations and World--Action Architectures}

% Recent world--action models use pretrained video representations to connect
% future-observation prediction with robot control. DreamZero, Cosmos3 and VideoVLA
% \cite{ye2026world, cosmos3, shen2026videovla} jointly generate future video and action
% chunks, while LingBot-VA~\cite{li2026causal} and DVA~\cite{rhoda2026dva} adopts a causal video-first formulation. EgoWAM and Zero-WAM~\cite{li2026egowam, zhou2026zero} use egocentric video as a source of predictive priors or task specifications, and ZimaBlue \cite{wu2026zimablue} studies multi-stage video pretraining followed by robot grounding.

% Closely related video-first approaches study how predictive visual information
% can be grounded into robot actions. Seer \cite{tian2025predictive} predicts actions
% through inverse dynamics conditioned on forecasted visual states, while VPP
% \cite{hu2024video} conditions an inverse-dynamics policy on predictive
% representations from a video model. AVDC \cite{ko2024learning} instead converts
% synthesized video trajectories into actions through dense visual
% correspondences. DreamGen \cite{jang2025dreamgen} uses video world models to
% generate synthetic robot trajectories and recovers pseudo-actions with
% inverse-dynamics or latent-action models. These methods provide complementary
% interfaces for translating predicted visual futures into executable behavior.

% A complementary line develops unified world--action models that expose multiple
% predictive and control modes within a shared architecture. UWM \cite{zhu2025unified} couples video and action diffusion and uses modality-specific diffusion timesteps to realize policy, forward-dynamics, inverse-dynamics, and video-generation modes. UVA \cite{li2503unified} jointly trains policy, inverse-dynamics, and forward-dynamics objectives within a shared
% video--action model. Riemann-1.0 and Motus2~\cite{sun2026riemann, bi2026motus2} similarly combine policy generation and action-conditioned prediction within broader shared-model systems; Motus2 additionally supports outcome evaluation and policy-improvement loops.

% Other work studies how predictive computation, representation, and supervision
% should be incorporated into world--action models. Fast-WAM, World Tokens, and
% GlanceWAM \cite{yuan2026fast, tang2026world, wang2026glancewam} reduce,
% sparsify, or reschedule future prediction during inference. SelfWAM
% \cite{pan2026selfwam} conditions future prediction on clean actions during
% training while retaining an action-only deployment path. OA-WAM and RepWAM
% \cite{liu2026oa, wang2026repwam} modify the predictive representation through
% object-addressable or action-grounded latents, while GIFT, SA-WAM, and TacPAC
% \cite{zheng2026gift, gonzalez2026spatially, ma2026tacpac} introduce structural,
% geometric, or tactile supervision. AcrossWAM1.0
% \cite{zhang2026acrosswam1} makes the boundaries between latent prediction,
% policy adaptation, and deployment explicit.

% Together, these works expose several coupled design choices, including the
% prediction target, cross-modal information flow, generation order,
% representation, supervision, and inference schedule. \ours fixes a common
% causal robot-video initialization and training interface while explicitly
% varying generation order and future cross-modal visibility, enabling joint,
% sequential, and decoupled video--action interaction to be compared under a
% shared evaluation protocol. Beyond within-model interaction, we also study
% whether independently trained dynamics components can be reused through
% explicit trajectory interfaces.

% \subsection{Video Foundations and World--Action Architectures}

% Recent world--action models use pretrained video representations to connect
% future prediction with robot control. DreamZero, Cosmos3, VideoVLA, and
% PAD~\cite{ye2026world, cosmos3, shen2026videovla, guo2024prediction}
% jointly model future visual content and actions, while LingBot-VA and
% DVA~\cite{li2026causal, rhoda2026dva} adopt causal video-first formulations.
% CoVAR and DeVA~\cite{yang2026covar, zhang2026deva} instead couple specialized
% video and action pathways through cross-modal feature exchange. Other work
% leverages egocentric or large-scale video priors for robot grounding and
% generalization~\cite{li2026egowam, zhou2026zero, wu2026zimablue,
% qiu2026vid2wam}.

% A closely related line studies video-to-action grounding. Seer and
% VPP~\cite{tian2025predictive, hu2024video} use predicted visual states or
% representations to condition inverse dynamics, AVDC~\cite{ko2024learning}
% converts synthesized videos to actions through dense correspondence, and
% DreamGen~\cite{jang2025dreamgen} recovers pseudo-actions for generated robot
% videos. DiT4DiT~\cite{ma2026dit4dit} couples action generation to intermediate
% video-denoising features, while World2Act~\cite{vuong2026world2act} transfers
% predictive dynamics through latent representation alignment.

% Unified formulations expose several predictive and control modes within one
% system. UWM and UVA~\cite{zhu2025unified, li2503unified} support combinations
% of policy, inverse-dynamics, forward-dynamics, and video prediction objectives.
% WorldVLA, Cosmos Policy, Riemann-1.0, Motus2, and X-WAM
% \cite{cen2025worldvla, kim2026cosmos, sun2026riemann, bi2026motus2,
% guo2026unified} further combine action generation with world prediction,
% simulation, planning, or evaluation in shared architectures.

% Recent work also isolates individual WAM design choices. Fast-WAM, World
% Tokens, GlanceWAM, GigaWorld-Policy, Light-WAM, Streaming-WAM, NoiseGate, and
% SelfWAM~\cite{yuan2026fast, tang2026world, wang2026glancewam,
% ye2603gigaworld, li2026light, huang2026streaming,
% huang2026noisegate, pan2026selfwam} study when, how much, and in what order
% predictive computation should interact with action generation. OA-WAM,
% RepWAM, AGRA, SG-WAM, LD4WAM, AffordanceWAM, and ReWAM
% \cite{liu2026oa, wang2026repwam, qiu2026making, zhao2026sg,
% shen2026ld4wam, you2026affordancewam, jiang2026rethinking} study alternative
% predictive representations, while GIFT, SA-WAM, and TacPAC
% \cite{zheng2026gift, gonzalez2026spatially, ma2026tacpac} introduce
% structural, geometric, or tactile supervision. WALL-WM, GE-Act 2.0,
% MotionWAM, and MaP-WAM~\cite{li2026wall, liu2026ge,
% zheng2026motionwam, zhao2026memory} extend this design space to temporal
% abstraction, scaling, humanoid control, and long-horizon memory; recent
% surveys summarize the rapidly expanding landscape
% \cite{wang2026world, lu2026world}.

% Concurrent work, also named OpenWAM~\cite{wang2026openwam}, develops a modular
% framework for systematic WAM pretraining and design studies. Our focus is
% complementary: \ours fixes a common causal robot-video foundation and training
% interface while explicitly varying generation order and future cross-modal
% visibility, enabling joint, video-then-action, action-then-video, and decoupled
% interaction to be compared under a shared evaluation protocol. We further study
% the reuse of independently trained inverse- and forward-dynamics components
% through explicit trajectory interfaces.

\subsection{Video Foundations and World--Action Architectures}

Recent world--action models connect predictive video representations with
robot control through several forms of video--action coupling. Joint
generation models predict future visual content and actions together
\cite{ye2026world, cosmos3, shen2026videovla, guo2024prediction}, while
causal video-first formulations generate or represent visual futures before
action prediction~\cite{li2026causal, rhoda2026dva}. Other architectures use
specialized video and action pathways with explicit cross-modal exchange
\cite{yang2026covar, zhang2026deva}, or leverage egocentric and large-scale
video priors for robot grounding and generalization
\cite{li2026egowam, zhou2026zero, wu2026zimablue, qiu2026vid2wam}.

A closely related line studies the interface between predictive visual
information and executable actions. Predicted visual states or representations
can condition inverse dynamics~\cite{tian2025predictive, hu2024video, li2026turningvideomodelsgeneralist}, while
other approaches translate generated futures through dense correspondence,
pseudo-action labeling, intermediate denoising features, or latent
representation alignment
\cite{ko2024learning, jang2025dreamgen, ma2026dit4dit,
vuong2026world2act}. These methods differ in where and how predictive visual
information enters the action-generation pathway.

Unified formulations instead expose multiple predictive and control modes
within a shared system. UWM and UVA~\cite{zhu2025unified, li2503unified}
support combinations of policy, inverse-dynamics, forward-dynamics, and video
prediction objectives, while related systems combine action generation with
world prediction, simulation, planning, or evaluation
\cite{cen2025worldvla, kim2026cosmos, sun2026riemann, bi2026motus2,
guo2026unified}.

Recent work further isolates particular choices in the WAM design space.
Predictive computation can be reduced, sparsified, reordered, or scheduled
asynchronously
\cite{yuan2026fast, tang2026world, wang2026glancewam,
ye2603gigaworld, li2026light, huang2026streaming,
huang2026noisegate, pan2026selfwam}. Alternative work studies the
representation through which world prediction interacts with control
\cite{liu2026oa, wang2026repwam, qiu2026making, zhao2026sg,
shen2026ld4wam, you2026affordancewam, jiang2026rethinking}, or augments
world--action learning with structural, geometric, and tactile supervision
\cite{zheng2026gift, gonzalez2026spatially, ma2026tacpac}. Temporal
abstraction, data and model scaling, humanoid control, and long-horizon memory
provide additional design axes
\cite{li2026wall, liu2026ge, zheng2026motionwam, zhao2026memory};
recent surveys summarize the rapidly expanding landscape
\cite{wang2026world, lu2026world}.

Concurrent work, also named OpenWAM~\cite{wang2026openwam}, develops a modular
framework for systematic WAM pretraining and design studies. Our focus is
complementary: \ours fixes a common causal robot-video foundation and training
interface while explicitly varying generation order and future cross-modal
visibility, enabling joint, video-then-action, action-then-video, and decoupled
interaction to be compared under a shared evaluation protocol. We further study
the reuse of independently trained inverse- and forward-dynamics components
through explicit trajectory interfaces.



% \subsection{Dynamics Learning Beyond Successful Demonstrations}

% Successful demonstrations provide valid transitions, but they cover only the
% action--outcome relationships selected by the demonstrating policy. A
% complementary analysis shows that observational demonstrations alone may not
% identify the consequences of unsupported actions~\cite{yu2026capability}.
% Several methods therefore use interactions beyond successful demonstrations. UniPi~\cite{du2023learning} separates video planning from inverse dynamics and allows the inverse model to be trained independently on a separate, potentially suboptimal dataset. PlayWorld~\cite{yin2026playworld} collects autonomous play for action-conditioned world modeling. FACT~\cite{peng2026fact} masks failed-action imitation while retaining action-conditioned future and progress supervision, and Motus2~\cite{bi2026motus2} routes failed, suboptimal, and task-irrelevant interactions to simulation or evaluation. WAM-OPD~\cite{yang2026wam} addresses a related distribution-shift problem by distilling a teacher on student-visited histories.
% % Several methods therefore exploit interactions beyond successful demonstrations. For example, UniPi~\cite{du2023learning} decouples video planning from inverse dynamics, allowing the inverse model to be trained on separate, potentially suboptimal data, while FACT~\cite{peng2026fact} retains future and progress supervision from failed trajectories without directly imitating their actions. Related work also explores autonomous play and other non-expert interactions for world modeling and evaluation.

% % Recent work has shown that failures, autonomous play, and counterfactual or
% % student-visited interactions can provide useful supervision for world modeling,
% % evaluation, or policy improvement. However, the effect of transition coverage
% % and conditioning scope on reusing independently trained dynamics components
% % after visual task adaptation remains unclear.
% Despite this progress, it remains unclear how transition coverage and conditioning scope affect the reuse of independently trained dynamics components after visual task adaptation.
% \ours addresses this question
% by evaluating local-context inverse- and forward-dynamics components and their
% composition with adapted video predictors under controlled changes in
% interaction data and conditioning context.
% \subsection{Dynamics Learning Beyond Successful Demonstrations}

% Successful demonstrations provide valid transitions, but they cover only the
% action--outcome relationships selected by the demonstrating policy. A
% complementary analysis shows that observational demonstrations alone may not
% identify the consequences of unsupported actions~\cite{yu2026capability}.
% Several methods therefore broaden supervision beyond clean expert
% demonstrations. DART~\cite{laskey2017dart} injects noise during expert
% demonstration collection to expose recovery states and reduce covariate shift
% in imitation learning. UniPi~\cite{du2023learning} separates video planning
% from inverse dynamics and allows the inverse model to be trained independently
% on a separate, potentially suboptimal dataset. PlayWorld~\cite{yin2026playworld}
% collects autonomous play for action-conditioned world modeling.
% FACT~\cite{peng2026fact} masks failed-action imitation while retaining
% action-conditioned future and progress supervision, and
% Motus2~\cite{bi2026motus2} routes failed, suboptimal, and task-irrelevant
% interactions to simulation or evaluation. WAM-OPD~\cite{yang2026wam}
% addresses a related distribution-shift problem by distilling a teacher on
% student-visited histories.

% \ours instead studies how broader local transition coverage interacts with
% conditioning scope to affect the reuse of independently trained dynamics
% components after visual task adaptation. To this end, we branch alternative
% action continuations from task-relevant simulator states and use the resulting
% state--action--outcome transitions as local dynamics supervision.

% \subsection{Dynamics Learning Beyond Successful Demonstrations}

% Successful demonstrations provide valid transitions, but they cover only the
% action--outcome relationships selected by the demonstrating policy. A
% complementary analysis shows that observational demonstrations alone may not
% identify the consequences of unsupported actions~\cite{yu2026capability}.
% A long line of imitation-learning work therefore broadens the training
% distribution beyond clean expert trajectories. DAgger~\cite{ross2011reduction}
% collects supervision on states visited by the learned policy, while
% DART~\cite{laskey2017dart} injects noise during expert demonstration collection
% to expose states arising from execution errors. Human-in-the-loop approaches
% such as IWR~\cite{mandlekar2020human} collect corrective interventions during
% policy rollouts, and more recent methods such as
% RaC~\cite{hu2026rac} explicitly collect recovery and correction trajectories
% for long-horizon manipulation. RECAP~\cite{intelligence2025pi} further combines
% demonstrations, autonomous on-policy experience, and expert interventions to
% improve vision--language--action policies through deployment.

% Recent work also constructs broader interaction data without relying solely on
% successful demonstrations. AHA~\cite{duan2024aha} procedurally
% perturbs successful simulated demonstrations to synthesize failure trajectories,
% while Dream2Fix~\cite{li2026learning} uses a generative world model to produce
% counterfactual failure rollouts from successful real-world demonstrations.
% UniPi~\cite{du2023learning} separates video planning from inverse dynamics and
% allows the inverse model to be trained independently on a separate, potentially
% suboptimal dataset. PlayWorld~\cite{yin2026playworld} collects autonomous play
% for action-conditioned world modeling. FACT~\cite{peng2026fact} masks
% failed-action imitation while retaining action-conditioned future and progress
% supervision, and Motus2~\cite{bi2026motus2} routes failed, suboptimal, and
% task-irrelevant interactions to simulation or evaluation.
% WAM-OPD~\cite{yang2026wam} addresses a related distribution-shift problem by
% distilling a teacher on student-visited histories.

% These approaches use broader experience primarily for robust imitation,
% recovery, policy improvement, failure understanding, or world modeling.
% \ours instead studies how local transition coverage interacts with conditioning
% scope to affect the reuse of independently trained dynamics components after
% visual task adaptation. To this end, we branch alternative action continuations
% from task-relevant simulator states and use the resulting
% state--action--outcome transitions as local dynamics supervision.

% \subsection{Dynamics Learning Beyond Successful Demonstrations}

% Successful demonstrations provide valid transitions, but they cover only the
% action--outcome relationships selected by the demonstrating policy. A
% complementary analysis shows that observational demonstrations alone may not
% identify the consequences of unsupported actions~\cite{yu2026capability}.
% A long line of imitation-learning work therefore broadens the training
% distribution beyond clean expert trajectories. DAgger~\cite{ross2011reduction}
% collects supervision on states visited by the learned policy, while
% DART~\cite{laskey2017dart} injects noise during expert demonstration collection
% to expose states arising from execution errors. Human-in-the-loop approaches
% such as IWR~\cite{mandlekar2020human} collect corrective interventions during
% policy rollouts, and more recent methods such as
% RaC~\cite{hu2026rac} explicitly collect recovery and correction trajectories
% for long-horizon manipulation. RECAP~\cite{intelligence2025pi} further combines
% demonstrations, autonomous on-policy experience, and expert interventions to
% improve vision--language--action policies through deployment.

% Recent work also constructs broader interaction data without relying solely on
% successful demonstrations. RoCoDA~\cite{ameperosa2025rocoda} synthesizes
% additional demonstrations through causal and geometric counterfactual
% augmentation, while AHA~\cite{duan2024aha} perturbs successful simulated
% demonstrations to generate failure trajectories. Dream2Fix
% \cite{li2026learning} uses a generative world model to synthesize
% counterfactual failure rollouts from successful real-world demonstrations,
% and FLARE~\cite{zhao2026flare} introduces perturbation-driven retry behavior
% and targeted reset skills for autonomous recovery. Fail2Progress
% \cite{huang2025fail2progress} generates data targeted to observed failures to
% improve skill-effect models. UniPi~\cite{du2023learning} separates video
% planning from inverse dynamics and allows the inverse model to be trained
% independently on a separate, potentially suboptimal dataset.
% PlayWorld~\cite{yin2026playworld} collects autonomous play for
% action-conditioned world modeling. FACT~\cite{peng2026fact} masks
% failed-action imitation while retaining action-conditioned future and progress
% supervision, and Motus2~\cite{bi2026motus2} routes failed, suboptimal, and
% task-irrelevant interactions to simulation or evaluation.
% WAM-OPD~\cite{yang2026wam} addresses a related distribution-shift problem by
% distilling a teacher on student-visited histories.

% These approaches use broader experience primarily for robust imitation,
% recovery, policy improvement, failure understanding, or world modeling.
% \ours instead studies how local transition coverage interacts with conditioning
% scope to affect the reuse of independently trained dynamics components after
% visual task adaptation. To this end, we branch alternative action continuations
% from task-relevant simulator states and use the resulting
% state--action--outcome transitions as local dynamics supervision.

\subsection{Dynamics Learning Beyond Successful Demonstrations}

Successful demonstrations provide valid transitions, but they cover only the
action--outcome relationships selected by the demonstrating policy. A
complementary analysis shows that observational demonstrations alone may not
identify the consequences of unsupported actions~\cite{yu2026capability}.
Classical imitation-learning methods address the resulting distribution shift
by collecting supervision on learner-visited states, perturbing expert
trajectories, or soliciting corrective interventions
\cite{ross2011reduction,laskey2017dart,mandlekar2020human,hu2026rac,
intelligence2025pi}. These approaches broaden state coverage beyond clean
expert rollouts and improve robustness to execution errors.

More recent work expands supervision through synthetic failures,
counterfactual augmentation, autonomous experience, and other non-successful
interactions
\cite{ameperosa2025rocoda,duan2024aha,li2026learning,zhao2026flare,
yin2026playworld,peng2026fact,bi2026motus2,yang2026wam}.
Such data have been used for recovery, policy improvement, failure
understanding, and predictive modeling. Particularly related to our setting,
Fail2Progress~\cite{huang2025fail2progress} targets data collection toward
observed failures to improve skill-effect models, while
UniPi~\cite{du2023learning} decouples video planning from inverse dynamics and
allows the inverse model to learn from a separate, potentially suboptimal
dataset.

\ours focuses on a different question: how transition coverage and conditioning
scope affect the reuse of independently trained dynamics components after
visual task adaptation. We branch alternative action continuations from
task-relevant simulator states and use the resulting state--action--outcome
transitions as local dynamics supervision.


% BEGIN ARCHIVED RELATED WORK -- superseded 2026-09-14
% \section{Related Work}
% \label{sec:related}
%
% \subsection{Visuomotor Policies and Predictive Control}
% Direct visuomotor policies, including ACT and Diffusion Policy
% \cite{zhao2023learning, chi2025diffusion}, learn executable action sequences from
% observations. RT-1/RT-2, OpenVLA, Octo, and the $\pi$ family
% \cite{brohan2022rt, brohan2023rt, kim2024openvla, team2024octo, black2024pi_0, intelligence2025pi_} extend this paradigm through
% language conditioning and large-scale pretraining, with Open X-Embodiment
% \cite{o2024open} supporting training across robot embodiments.
% Predictive approaches additionally model future observations or latent states:
% Dreamer and DayDreamer \cite{hafner2019dream, hafner2023mastering, wu2023daydreamer} learn dynamics for
% control, while GR-1/GR-2 and RoboDreamer \cite{wu2024unleashing, cheang2024gr, zhou2024robodreamer}
% incorporate future prediction into policy learning or planning. WorldEval and
% Ctrl-World \cite{li2025worldeval, guo2026ctrl} use predicted rollouts for policy
% assessment, and A2World and $\tau_0$-WM \cite{huang2026learning, zhou2026tau_0} develop
% action-conditioned models for simulation and control. World-Coherent Decoding
% \cite{zhang2026worldcoherent} further uses predicted future--action pairs for
% selection and execution-time verification. These approaches motivate
% predictive interfaces with distinct roles in action generation, simulation,
% and evaluation.
%
% \subsection{Video Foundations and World--Action Architectures}
% Pretrained video generators provide visual and temporal priors for robot control. DreamZero~\cite{ye2026world} and VideoVLA~\cite{shen2026videovla} jointly
% generate future video and actions, while LingBot-VA~\cite{li2026causal}
% adopts a causal video-first formulation. EgoWAM~\cite{li2026egowam}, Zero-WAM~\cite{zhou2026zero}, and ZimaBlue~\cite{wu2026zimablue} explore egocentric video as a source of predictive priors or task specifications. \ours builds on this direction through an explicit stage of causal robot-video adaptation without action labels, providing a common initialization for downstream video--action learning.
%
% Recent WAMs explore several additional design dimensions. Fast-WAM, World
% Tokens, RIFT, and GlanceWAM
% \cite{yuan2026fast, tang2026world, zhang2026keep, wang2026glancewam} reduce or
% reschedule predictive computation during control. OA-WAM, RepWAM, GIFT,
% SA-WAM, and TacPAC \cite{liu2026oa, wang2026repwam, zheng2026gift, gonzalez2026spatially, ma2026tacpac}
% investigate alternative representations, structural supervision, or tactile
% prediction. Riemann-1.0 and Motus2 \cite{sun2026riemann, bi2026motus2} support
% policy generation and action-conditioned simulation through shared parameters,
% with Motus2 also providing outcome evaluation. Differences in backbones,
% supervision, and inference procedures complicate comparisons across systems.
% \ours focuses on controlled comparisons of joint, sequential, and restricted
% video--action interaction under a common architecture, examining both policy
% performance and generation cost.
%
% \subsection{Dynamics Learning Beyond Successful Demonstrations}
% Successful demonstrations provide useful transitions but cover only a subset
% of possible action--outcome relationships. A complementary analysis shows why
% observational demonstrations alone need not identify unsupported action
% effects \cite{yu2026capability}. Learning from other interactions has
% established precedents. UniPi \cite{du2023learning} separates video planning from
% inverse dynamics and explicitly allows the inverse model to be trained
% independently on a separate, potentially suboptimal dataset. PlayWorld
% \cite{yin2026playworld} collects autonomous play for world modeling. FACT
% \cite{peng2026fact} uses failures to supervise action-conditioned video and
% progress prediction while masking failed-action imitation; Motus2
% \cite{bi2026motus2} routes failed, suboptimal, and task-irrelevant interactions
% to simulation or evaluation. WAM-OPD \cite{yang2026wam} addresses a
% complementary source of distribution shift by distilling a teacher on
% student-visited histories.
%
% Our study builds on these directions by examining transition coverage and
% conditioning scope as separate factors in dynamics-component reuse. We
% independently train inverse- and forward-dynamics components on local
% transitions, using demonstrations and counterfactual rollouts in simulation,
% and extending inverse-dynamics supervision with physical-robot failures and
% play. Counterfactual collection executes alternative action sequences from the
% same simulator state; each resulting transition supplies an individual
% dynamics-training example. The strict interfaces omit task language and
% preceding history. Comparisons with full-context models under matched
% supervision examine this restriction separately from the effect of broader
% data. Component-composition experiments then evaluate whether frozen
% source-trained inverse models remain useful when paired with task-adapted
% video predictors.
%
% \subsection{Open Frameworks for Controlled WAM Studies}
% OpenVLA, Octo, StarVLA, VLA Foundry, and AcrossWAM
% \cite{kim2024openvla, team2024octo, ye2026starvla, mercat2026vla, zhang2026acrosswam1}
% motivate reusable infrastructure for embodied learning. \ours contributes a
% common robot-video foundation and configurable video--action architecture,
% together with shared training, rollout, and evaluation interfaces. These
% interfaces support both architecture comparisons and composition of
% independently trained components, connecting studies of video--action
% interaction with studies of dynamics supervision and reuse.
% END ARCHIVED RELATED WORK


% BEGIN ARCHIVED RELATED WORK -- superseded 2026-09-10
% \section{Related Work}
% \label{sec:related}
%
% \subsection{Direct Visuomotor and Vision-Language-Action Policies}
% Direct policies map observations and, increasingly, language instructions to
% robot actions without an explicit predictive world interface. RT-1/RT-2
% \cite{brohan2022rt, brohan2023rt} established large-scale language-conditioned control, Open
% X-Embodiment and Octo \cite{o2024open, team2024octo} scaled training across robot
% embodiments, and OpenVLA and the $\pi$ family
% \cite{kim2024openvla, black2024pi_0, intelligence2025pi_} transferred pretrained vision-language features
% to generalist manipulation. ACT and Diffusion Policy
% \cite{zhao2023learning, chi2025diffusion} further show that action chunking and generative
% action heads are themselves important design choices. These methods provide
% strong direct-control baselines: they may learn dynamics implicitly from
% history and demonstrations, but their objectives do not require a separately
% addressable prediction of future scene evolution.
%
% \subsection{Video-Predictive Policies and World Models}
% Predictive models have long supported planning, representation learning, and
% simulation. Dreamer and DayDreamer \cite{hafner2019dream, hafner2023mastering, wu2023daydreamer} learn
% latent dynamics for control, while UniPi \cite{du2023learning} generates video plans
% before recovering actions. GR-1/GR-2 and RoboDreamer
% \cite{wu2024unleashing, cheang2024gr, zhou2024robodreamer} use future prediction as policy supervision or a
% compositional plan. Other work treats prediction as a simulator or evaluator:
% WorldEval and Ctrl-World \cite{li2025worldeval, guo2026ctrl} roll out controlled
% futures for policy assessment, while A2World and $\tau_0$-WM
% \cite{huang2026learning, zhou2026tau_0} learn action-conditioned dynamics that can serve
% both simulation and control. World-Coherent Decoding
% \cite{zhang2026worldcoherent} further treats predicted future--action pairs as
% hypotheses to select and verify after execution. These approaches establish the
% value of predictive supervision, but connect it to control through different
% planners, inverse models, evaluators, and training pipelines.
%
% \subsection{World-Action Model Architectures}
% Recent WAMs make video--action coupling part of the policy itself. DreamZero
% \cite{ye2026world} and VideoVLA \cite{shen2026videovla} jointly predict future
% video and action chunks using pretrained video generators, whereas LingBot-VA
% \cite{li2026causal} follows a causal video-first formulation. A rapidly
% growing line reduces the runtime cost of imagination: Fast-WAM, World Tokens,
% and RIFT \cite{yuan2026fast, tang2026world, zhang2026keep} retain predictive
% training or anticipation while omitting dense video rollout from the action
% path, and GlanceWAM \cite{wang2026glancewam} schedules sparse lookahead
% asynchronously. Another line changes what the predictive interface represents:
% OA-WAM and RepWAM \cite{liu2026oa, wang2026repwam} learn object-addressable or
% action-grounded latents, GIFT and SA-WAM \cite{zheng2026gift, gonzalez2026spatially} inject
% structural or 3D supervision, and TacPAC \cite{ma2026tacpac} adds tactile
% prediction and correction. Riemann-1.0 and Motus2
% \cite{sun2026riemann, bi2026motus2} expose policy and action-conditioned simulation
% through shared weights; Motus2 additionally exposes outcome evaluation. Since
% these systems vary backbone, target, coupling, data, and inference schedule
% together, their results do not isolate which design choice causes a gain.
% Open-WAM instead makes these choices interchangeable under a shared causal
% video backbone and common training, rollout, and evaluation contracts.
%
% \subsection{Data Utilization and Open WAM Infrastructure}
% Most manipulation WAMs are post-trained primarily on successful task
% demonstrations, even though broader interaction data also contains useful
% physical transitions. Several recent works directly address this limitation.
% PlayWorld \cite{yin2026playworld} collects autonomous play for contact-rich world
% modeling; FACT \cite{peng2026fact} masks imitation loss on failed actions while
% retaining their action-conditioned video and progress supervision; and Motus2
% routes failed, suboptimal, and task-irrelevant interactions to simulation and
% evaluation modes. WAM-OPD \cite{yang2026wam} instead distills a teacher on
% student-visited histories. A complementary analysis shows why observational
% demonstrations alone need not identify unsupported action effects
% \cite{yu2026capability}. Human-video scaling provides another source
% of broader physical experience: EgoWAM, Zero-WAM, and ZimaBlue
% \cite{li2026egowam, zhou2026zero, wu2026zimablue} learn predictive priors or task
% specifications from egocentric video. These advances narrow the data gap but
% leave a heterogeneous set of task-specific simulators, value heads, or
% post-training loops. \ours-GJD takes a complementary route: a single shared
% vector field switches among task-conditioned joint policy generation and
% strict, task-agnostic IDM/FDM modes, so successful demonstrations and paired
% counterfactual local transitions can supervise the same model without treating
% undesirable actions as policy targets. In parallel, OpenVLA, Octo, StarVLA, VLA
% Foundry, and AcrossWAM
% \cite{kim2024openvla, team2024octo, ye2026starvla, mercat2026vla, zhang2026acrosswam1}
% motivate reusable training and deployment stacks; Open-WAM extends this
% principle to controlled WAM architecture and data-utilization studies.
% END ARCHIVED RELATED WORK
